\documentclass[a4paper,12pt]{article}
\usepackage[utf8]{inputenc}
\usepackage[T1]{fontenc}
\usepackage[ngerman]{babel}
\usepackage{amsmath, amssymb}
\usepackage{geometry}
\usepackage{parskip}
\usepackage{titlesec}
\usepackage{enumitem}

\geometry{left=3cm, right=3cm, top=2.5cm, bottom=2.5cm}
\titleformat{\section}{\large\bfseries}{}{0em}{}[\titlerule]

\newcommand{\gentry}[3]{%
  \noindent\textbf{#1\quad #2}\nopagebreak\\[0.2em]
  #3\par\smallskip
}

\begin{document}

\begin{center}
  {\LARGE\bfseries Glossar zu Lektion 5}\\[0.5em]
  {\large Lineare Algebra (Modul 61112)}
\end{center}

\vspace{0.8em}

% ======================================================================
\section*{5.1 \quad Normalformproblem}

\gentry{5.1.2}{Normalformproblem}{%
  Zu $\varphi \in \mathrm{End}(V)$: Finde eine geordnete Basis $B$, sodass $M_B(\varphi)$ möglichst einfach ist.
  Äquivalent für Matrizen: Finde ein Vertretersystem der Ähnlichkeitsklassen in $M_{n,n}(K)$.
  Zwei Endomorphismen sind \textbf{ähnlich} (bzgl.\ Basiswechsel), wenn ihre Matrixdarstellungen ähnlich sind: $B' = S^{-1}BS$.
}

\gentry{5.1.5}{Invarianter Unterraum}{%
  $U \subseteq V$ heißt $\varphi$-\textbf{invariant}, wenn $\varphi(u) \in U$ für alle $u \in U$.
  Dann ist $\varphi|_U : U \to U$ ein Endomorphismus von $U$.
}

\gentry{5.1.13}{Satz (Blockdiagonalform)}{%
  Ist $V = U_1 \oplus \cdots \oplus U_k$ mit $\varphi$-invarianten Unterräumen, so hat $M_B(\varphi)$ bzgl.\ der vereinigten Basis Blockdiagonalgestalt.
}

% ======================================================================
\section*{5.2 \quad Eigenwerte, Eigenvektoren, Diagonalisierbarkeit}

\gentry{5.2.1}{Eigenvektor, Eigenwert}{%
  $v \neq 0$ heißt \textbf{Eigenvektor} von $\varphi$ zum \textbf{Eigenwert} $\lambda \in K$, wenn $\varphi(v) = \lambda v$.
  $\lambda$ ist Eigenwert $\iff$ $\varphi - \lambda \cdot \mathrm{id}_V$ ist \emph{nicht} bijektiv.
}

\gentry{5.2.8}{Eigenraum}{%
  Zu $\lambda \in K$: $E_\lambda(\varphi) := \ker(\varphi - \lambda \cdot \mathrm{id}_V)$. Dies ist ein $\varphi$-invarianter Unterraum.
}

\gentry{5.2.10}{Geometrische Vielfachheit}{%
  $\mathrm{gVfh}(\lambda) := \dim E_\lambda(\varphi) \geq 1$ (wenn $\lambda$ Eigenwert).
}

\gentry{5.2.13}{Lemma (Summe von Eigenräumen ist direkt)}{%
  Die Summe der Eigenräume zu verschiedenen Eigenwerten ist direkt:
  $E_{\lambda_1} + \cdots + E_{\lambda_r} = E_{\lambda_1} \oplus \cdots \oplus E_{\lambda_r}$.
}

\gentry{5.2}{Diagonalisierbarkeit}{%
  $\varphi$ heißt \textbf{diagonalisierbar}, wenn $V$ eine Basis aus Eigenvektoren besitzt.
  Äquivalent: $V = E_{\lambda_1} \oplus \cdots \oplus E_{\lambda_r}$ (direkte Summe aller Eigenräume).
}

% ======================================================================
\section*{5.3 \quad Charakteristisches Polynom und Minimalpolynom}

\gentry{5.3.1}{Charakteristisches Polynom}{%
  $\chi_\varphi(X) := \det(X \cdot \mathrm{id}_V - \varphi) \in K[X]$. Dies ist ein normiertes Polynom vom Grad $n = \dim V$.
  \textbf{Eigenwerte} von $\varphi$ sind genau die Nullstellen von $\chi_\varphi$.
}

\gentry{5.3}{Algebraische Vielfachheit}{%
  Die \textbf{algebraische Vielfachheit} $\mathrm{aVfh}(\lambda)$ ist die Vielfachheit von $\lambda$ als Nullstelle von $\chi_\varphi$.
  Es gilt stets: $\mathrm{gVfh}(\lambda) \leq \mathrm{aVfh}(\lambda)$.
}

\gentry{5.3.32}{Minimalpolynom}{%
  Das \textbf{Verschwindungsideal} $I(\varphi) := \{f \in K[X] \mid f(\varphi) = 0\}$ ist ein Ideal in $K[X]$.
  Der normierte Erzeuger $\mu_\varphi$ heißt \textbf{Minimalpolynom} von $\varphi$.
  Es ist der eindeutige normierte Teiler kleinsten Grades mit $\mu_\varphi(\varphi) = 0$.
}

\gentry{5.3.39}{Satz (Nullstellen des Minimalpolynoms)}{%
  Die Nullstellen von $\mu_\varphi$ sind genau die Eigenwerte von $\varphi$.
}

\gentry{5.3.42}{Satz von Cayley-Hamilton}{%
  $\chi_\varphi(\varphi) = 0$, d.\,h.\ $\mu_\varphi \mid \chi_\varphi$.
}

% ======================================================================
\section*{5.4 \quad Stochastische Matrizen}

\gentry{5.4.1}{Stochastische Matrix}{%
  $A \in M_{n,n}(\mathbb{R})$ heißt \textbf{(spalten)stochastisch}, wenn alle Einträge $\geq 0$ sind und jede Spalte die Summe~$1$ hat: $\sum_i a_{ij} = 1$.
}

\gentry{5.4.6}{Satz von Perron}{%
  Jede \textbf{positive} stochastische Matrix ($a_{ij} > 0$) hat den Eigenwert~$1$ mit geometrischer Vielfachheit~$1$. Alle anderen Eigenwerte haben Betrag $< 1$.
}

\end{document}
